% META: {"id": "TNSI-STRUCT-EX-020", "chapitre": "TNSI-STRUCTURES-DONNEES", "type_objet": "exercice", "capacites": ["TNSI-STRUCTURES-DONNEES-C5"], "capacites_codes": ["C5"], "methodes": ["M1"], "parcours": 2, "competences": ["programmer", "analyser"], "duree_min": 8, "mode_creation": "ex_nihilo", "sources_inspiration": [], "fichier_tex": "chapitres/TNSI-STRUCTURES-DONNEES/exercices/TNSI-STRUCT-EX-020.tex", "corrige_tex": "chapitres/TNSI-STRUCTURES-DONNEES/corriges/TNSI-STRUCT-CO-020.tex", "status": "approved"}
% BEGIN-VERIFY
% matrice = [[0,1,0,1],[0,0,1,0],[0,0,0,0],[0,0,1,0]]
% succ = {i: [j for j in range(4) if matrice[i][j]==1] for i in range(4)}
% assert succ == {0: [1,3], 1: [2], 2: [], 3: [2]}
% END-VERIFY
\begin{exercice}{TNSI-STRUCT-EX-004}{2}{18}
On donne le graphe oriente a $4$ sommets ($0,1,2,3$) represente par la matrice d'adjacence :
\begin{python}
matrice = [
    [0, 1, 0, 1],
    [0, 0, 1, 0],
    [0, 0, 0, 0],
    [0, 0, 1, 0],
]
\end{python}
\begin{enumerate}
  \item Lister toutes les aretes de ce graphe (sous la forme $i \to j$).
  \item Le sommet $2$ a-t-il des successeurs ? Interpreter.
  \item En utilisant la fonction \lstinline{matrice_vers_listes} du cours, donner la representation par listes de successeurs de ce graphe.
\end{enumerate}
\end{exercice}
